%!TEX root = ../thesis.tex

\chapter{Instructions to Compile and Run System}
\label{apx:instructions}

Write your Appendix content here.

\section{Section Heading the Appendix}
\label{apx:instructions:section}

This is how you can write the content of your Appendix.
You can have multiple sections and subsections as needed.

\begin{instructions}
\section{Template Resources}
\label{apx:instructions:resources}

\begin{itemize}
\item Overleaf template: \url{https://www.overleaf.com/latex/templates/uis-thesis-template/qhfrhftzwqbc}
\item Thesis writing at the Faculty of Science and Technology (English): \url{https://www.uis.no/en/student-pages/studies/thesis-writing-at-faculty-of-science-and-technology}
\item Oppgaveskriving ved det teknisk-naturvitenskapelige fakultet (norsk): \url{https://www.uis.no/nb/studentsider/studier/oppgaveskriving-ved-det-teknisk-naturvitenskapelige-fakultet}
\end{itemize}

\section{Examples of Commonly Used Packages}
\label{apx:instructions:packages}

The packages below are loaded in \texttt{thesis.tex}.
Remove those that you do not need.

\paragraph{cleveref.}
\cref{apx:instructions} and \cref{apx:instructions:packages} are cross-referenced with \verb|\cref{label}|; the type (chapter, section, \dots) is added automatically.

\paragraph{booktabs and multirow.}
\Cref{tab:example} shows a table with \texttt{booktabs} rules and a \texttt{multirow} cell.
\begin{table}[t]
\centering
\caption{Example table.}
\label{tab:example}
\begin{tabular}{llr}
\toprule
Group & Item & Value \\
\midrule
\multirow{2}{*}{A} & first  & 1 \\
                   & second & 2 \\
B                  & third  & 3 \\
\bottomrule
\end{tabular}
\end{table}

\paragraph{algorithm and algpseudocode.}
\Cref{alg:example} shows a floating algorithm.
\begin{algorithm}[t]
\caption{Sum of a list}
\label{alg:example}
\begin{algorithmic}[1]
\Procedure{Sum}{$a_1,\dots,a_n$}
  \State $s \gets 0$
  \For{$i \gets 1$ \textbf{to} $n$}
    \State $s \gets s + a_i$
  \EndFor
  \State \Return $s$
\EndProcedure
\end{algorithmic}
\end{algorithm}

\paragraph{listings.}
\Cref{lst:example} shows a source code listing.
\begin{lstlisting}[caption={Example listing.},label={lst:example},basicstyle=\ttfamily\small]
func Sum(a []int) (s int) {
	for _, v := range a {
		s += v
	}
	return
}
\end{lstlisting}

\paragraph{xspace.}
Define macros such as \verb|\newcommand{\sys}{MySystem\xspace}| so that spacing after the macro is handled correctly.

\paragraph{Other useful packages.}
Consider \texttt{pgfplots} and \texttt{pgfplotstable} for plots and tables from data, and \texttt{acronym} for managing abbreviations.
\end{instructions}
